\section{Examples of Boundedness and Unboundedness}\label{sec:examples}

This section reports running examples of smart contracts
and boundedness for containers. The examples are written in a pseudocode
inspired by Solidity and Takamaka~\cite{OlivieriST25}
(the latter provides a clearer type system, inherited from Java),
in order to concentrate on the semantics of the code rather than
on the idiosyncrasies of a specific programming language.

\begin{figure}[t]
  \begin{lstlisting}[language=pseudo]
public class Ponzi1 extends Contract {
  private final List<Contract> investors;|\label{ponzi1:investors}|
  private final static int MINIMUM_INVESTMENT = 1000;

  public Ponzi1() {
    investors = new List();
    investors.push(caller);|\label{ponzi1:creator}|
  }

  public payable void invest() {|\label{ponzi1:invest}|
    require(amount >= MINIMUM_INVESTMENT, "you invested too little");|\label{ponzi1:require}|

    int eachInvestorGets = amount / investors.size();
    for investor in investors|\label{ponzi1:iteration}|
      investor.receive(eachInvestorGets);|\label{ponzi1:receive}|

    investors.push(caller);|\label{ponzi1:expand}|
  }
}
  \end{lstlisting}
  \caption{A smart contract implementing a Ponzi scheme. Investors get their money
    back only if new investors keep arriving after them. The list \<investors>
    is unbounded and, consequently, the loop at line~\ref{ponzi1:iteration} is unbounded as well.}
  \label{fig:ponzi1}
\end{figure}

The smart contract in Fig.~\ref{fig:ponzi1}, from~\cite{IyerD18}, implements a Ponzi scheme.
It keeps a list of \<investors> (more precisely, \<this.investors>) at
line~\ref{ponzi1:investors} where, as in Takamaka, \<Contract> is the superclass of both
contracts and externally-owned accounts (treated as a special case of contracts).
The caller of the constructor populates that list initially, for free
(line~\ref{ponzi1:creator}, where \<caller> corresponds to
\<msg.sender> in Solidity). Subsequently, new contracts become investors
by calling method \<invest()> and paying a minimal fee (\<amount> corresponds to
\<msg.value> in Solidity). This fee gets split among the existing investors
(lines~\ref{ponzi1:iteration} and~\ref{ponzi1:receive}) before the \<caller> is
added to \<investors> (line~\ref{ponzi1:expand}). There is no upper bound
to the size of \<investors> and no way of removing investors.
As long as \<invest()> gets called,
\<investors> grows arbitrarily. Therefore, there is no upper bound
to the gas needed to execute the loop at line~\ref{ponzi1:iteration}: once
it exceeds the maximum allowed for a blockchain transaction, the smart contract is
effectively blocked: \<invest()> can only fail for gas exhaustion, reverting
all transaction updates, and some investors will never be able to redeem their investment.
We will say that list \<investors> and the loop at line~\ref{ponzi1:iteration}
are \emph{unbounded} and we expect a sound boundedness analyzer to warn about an
unbouded loop at line~\ref{ponzi1:iteration}\footnote{The \<receive()> at
line~\ref{ponzi1:receive} in Fig.~\ref{fig:ponzi1} is the payment
method that sends cryptocurrency to a \<Contract>. If it can be redefined
(as in Solidity but differently from Takamaka, where it is \<final>), it becomes
possible to create a re-entrancy attack. This is irrelevant here and
we won't discuss it further: the solution of using the \emph{withdrawal pattern}
still leaves this contract suffering from an unboundedness issue.}.
%
\begin{figure}[t]
  \begin{lstlisting}[language=pseudo]
public class Ponzi2 extends Contract {
  private final List<Contract> investors;|\label{ponzi2:investors}|
  private final static int MINIMUM_INVESTMENT = 1000;

  public Ponzi2() {
    investors = new List();
    investors.push(caller);|\label{ponzi2:creator}|
  }

  public payable void invest() {|\label{ponzi2:invest}|
    require(amount >= MINIMUM_INVESTMENT, "you invested too little");|\label{ponzi2:require}|
    require(investors.size() < 100, "no more investors are allowed");|\label{ponzi2:require_limit}|

    int eachInvestorGets = amount / investors.size();
    for investor in investors|\label{ponzi2:iteration}|
      investor.receive(eachInvestorGets);|\label{ponzi2:receive}|

    investors.push(caller);|\label{ponzi2:expand}|
  }
}
  \end{lstlisting}
  \caption{A modification of the contract in Fig.~\ref{fig:ponzi1}, that enforces
    an upper bound to the size of \<investors> (line~\ref{ponzi2:require_limit}).
    That list and the loop at line~\ref{ponzi2:iteration} are bounded.}
  \label{fig:ponzi2}
\end{figure}
%
Fig.~\ref{fig:ponzi2} shows how to fix this unboundedness issue.
It reports a modification of the smart contract in Fig.~\ref{fig:ponzi1} that
constrains the size of \<investors> (line~\ref{ponzi2:require_limit}) to be
$100$ at most. Thus, the iteration at line~\ref{ponzi2:iteration}
is bounded and a precise analyzer should issue no warning in Fig.~\ref{fig:ponzi2}.
This example shows that boundedness is unrelated to the exact maximal gas cost of a loop:
$100$ iterations might be acceptable or not, depending on the maximal gas allowed
for a blockchain transaction. Boundedness is an abstract semantical and blockchain-agnostic
property about the definitely limited size of iterated containers.

\begin{figure}
  \begin{lstlisting}[language=pseudo]
public class Lottery1 extends Contract {
  public final int TICKET_PRICE = 1000;|\label{lottery1:price}|
  private final List<Round> rounds; // all existing rounds|\label{lottery1:rounds}|
  private final long duration; // round duration|\label{lottery1:duration}|

  private static class TicketBatchSale {|\label{lottery1:sale}|
    private final Contract buyer;|\label{lottery1:buyer}|
    private final int tickets;

    private TicketBatchSale(Contract buyer, int tickets) {|\label{lottery1:constructor_tbs}|
      this.buyer = buyer;
      this.tickets = tickets;
    }
  }

  private static class Round { // ticket sale round|\label{lottery1:round}|
    private final long endTime;
    private final List<TicketBatchSale> sales;|\label{lottery1:sales}|
    private int totalTickets;|\label{lottery1:total_tickets}|
    private Contract winner;|\label{lottery1:winner}|

    private Round(long endTime) {
      this.endTime = endTime;
      sales = new List();
    }
  }

  public Lottery1(long duration) {
    this.duration = duration;
    rounds = new List();
    rounds.push(new Round(now + duration));|\label{lottery1:first_round}|
  }

  public payable void buy() {
    require(amount > 0 && amount % TICKET_PRICE == 0);|\label{lottery1:amount}|

    Round last = rounds.last();|\label{lottery1:last_start}|
    if (now > last.endTime) {
      last = new Round(now + duration);
      rounds.push(last);|\label{lottery1:expand}|
    }|\label{lottery1:last_end}|

    int quantity = amount / TICKET_PRICE;
    last.sales.push(new TicketBatchSale(caller, quantity));|\label{lottery1:add_sale}|
    last.totalTickets += quantity;
  }

  public void drawWinner(int roundNumber) {
    Round round = rounds[roundNumber];
    require(round.winner == NULL, "winner already drawn");|\label{lottery1:no_winner_yet}|
    require(now > round.endTime, "too early");|\label{lottery1:expired}|
    require(round.totalTickets > 0, "no tickets in round");|\label{lottery1:at_least_one_ticket}|

    int numWinningTicket = random(round.totalTickets);|\label{lottery1:random}|
    for sale in round.sales|\label{lottery1:iteration}|
      if (sale.tickets > numWinningTicket) {
        round.winner = sale.buyer;|\label{lottery1:set_winner}|
        round.winner.receive(TICKET_PRICE * round.totalTickets);|\label{lottery1:win}|
        return;|\label{lottery1:return}|
      }
      else
        numWinningTicket -= sale.tickets;
  }
}
  \end{lstlisting}
  \caption{A smart contract for a recurring lottery. Lottery rounds have a duration,
    after which a round expires. Tickets are sold in batches. This contract
    suffers from an unbounded iteration at line~\ref{lottery1:iteration} on the unbounded
    list of sales in a round.}
  \label{fig:lottery1}
\end{figure}

Fig.~\ref{fig:lottery1} shows another example from~\cite{IyerD18}, for a recurrent lottery.
Tickets have a fixed price (line~\ref{lottery1:price}) and are bought in batches
(inner class \<TicketBatchSale> at line~\ref{lottery1:sale}). The contract
keeps an unbounded list of \<rounds> (line~\ref{lottery1:rounds}), each lasting
for \<duration> milliseconds (line~\ref{lottery1:duration}) and containing
the list of batch ticket sales (line~\ref{lottery1:sales})
and the total number of tickets sold in the round (line~\ref{lottery1:total_tickets}).
Each round gets assigned a winner eventually (line~\ref{lottery1:winner})\ie one
of the buyers of its batch sales (line~\ref{lottery1:buyer}). The constructor
initializes \<rounds> with a single round, that expires after
\<duration> milliseconds from now (line~\ref{lottery1:first_round}).
Method \<buy()> is called by who wants to buy a batch of tickets. It requires
to pay enough cryptocurrency to buy exactly a batch of a
positive number of tickets (line~\ref{lottery1:amount}).
These tickets are bought for the last, most recent lottery round,
ensuring that a fresh new last round is added to \<rounds> if its last
round is already expired (lines~\ref{lottery1:last_start}--\ref{lottery1:last_end}).
Method \<drawWinner()> draws the winner of a round, if this has not been drawn yet
(line~\ref{lottery1:no_winner_yet}), if the round has already expired
(line~\ref{lottery1:expired}) and if at least one ticket has been sold in the round
(line~\ref{lottery1:at_least_one_ticket}). In such a case, the number of
the winning ticket is randomly chosen between 0 (included) and the total number of
tickets in the round (excluded) (line~\ref{lottery1:random}).
Then the sales in the round are iterated (line~\ref{lottery1:iteration}), looking for the
buyer of the winning ticket, which gets assigned as the winner of the round
(line~\ref{lottery1:set_winner}) and receives all money collected from the ticket sales
in the round (line~\ref{lottery1:win}).
Method \<drawWinner()> is meant to be called by any actor, most likely by somebody
who bought some ticket in the round.
The iteration at line~\ref{lottery1:iteration} is unbounded, since there is
no upper bound to the number of sales in a round (line~\ref{lottery1:add_sale}).
Therefore, the contract in Fig.~\ref{fig:lottery1}
suffers from an unboundedness bug and a sound analyzer should warn
at line~\ref{lottery1:iteration} about that.
In practice, this means that there is a limit $X$
to the number of sales that that iteration can scan before running out of gas.
Therefore, a sale has a chance to win the round only if it is among the first $X$
sales of the round (thanks to the \<return> at line~\ref{lottery1:return}, that
exits the loop). Therefore,
an attacker might flood a new round with $X$ one-ticket sales and be sure to win
the round, by repeatedly calling \<drawWinner()> until a
$\<numWinningTicket>\le X$ is drawn. Other players could be unaware of the bug and
subsequently buy extra tickets without realizing that they have no chance to win and,
by buying tickets, are only increasing the jackpot finally collected by the attacker.
Note that also \<rounds> (line~\ref{lottery1:rounds}) is unbounded,
but it is not iterated,
hence its unboundedness is not problematic and a precise analyzer is expected to
issue no warning about \<rounds>.
%
\begin{figure}[t]
  \begin{lstlisting}[language=pseudo,numbers=none]
public class Lottery2 extends Contract {
  ...
  private final List<Round> rounds; // all existing rounds
  public final static int MAX_SALES = 5000;|\label{lottery2:max_sales}|
  ...
  public payable void buy() {
    require(amount > 0 && amount % TICKET_PRICE == 0);

    Round last = rounds.last();
    if (now > last.endTime) {
      last = new Round(now + duration);
      rounds.push(last);|\label{lottery2:push_round}|
    }

    if (last.sales.size() < MAX_SALES) {
      int quantity = amount / TICKET_PRICE;
      last.sales.push(new TicketBatchSale(caller, quantity));
      last.totalTickets += quantity;
    }
  }
  ...
}
  \end{lstlisting}
  \caption{A modification of the smart contract in Fig.~\ref{fig:lottery1},
    where the list of sales in a round is limited. This version does not
    suffer from any unboundedness bug.}
  \label{fig:lottery2}
\end{figure}
%
A solution to this unboundedness bug is to
limit the maximal number of sales in each round, as shown in
Fig.~\ref{fig:lottery2}. Note that the list \<rounds> remains unbounded, but that
is irrelevant. We expect a precise boundedness analyzer to issue no
warning in Fig.~\ref{fig:lottery2}.
